\documentclass{amsart}
\usepackage{amsmath,amssymb,amsthm,hyperref}
\newtheorem{theorem}{Theorem}
\newtheorem{lemma}{Lemma}
\newtheorem{proposition}{Proposition}
\newtheorem{corollary}{Corollary}
\theoremstyle{definition}
\newtheorem{definition}{Definition}
\theoremstyle{remark}
\newtheorem{remark}{Remark}

\begin{document}

\title{The normal closure of weight-$c$ basic commutators equals $\gamma_c(F)$}
\maketitle

\section{Statement and conventions}

Let $F$ be a free group of finite rank $r$ with fixed free generating set
$X=\{x_1,\dots,x_r\}$. We write $X^{\pm}=X\cup X^{-1}$. We use the commutator
convention
\[
[a,b]=a^{-1}b^{-1}ab ,\qquad a^{b}=b^{-1}ab .
\]
The lower central series is $\gamma_1(F)=F$ and $\gamma_{k+1}(F)=[\gamma_k(F),F]$
for $k\ge 1$; for subgroups $H,K\le F$,
\[
[H,K]=\big\langle\, [h,k] : h\in H,\ k\in K \,\big\rangle .
\]
Each $\gamma_k(F)$ is fully invariant, hence normal, and
$\gamma_{k+1}(F)\le\gamma_k(F)$.

Basic commutators in $X$ and their weights are defined by the Hall ordering as in
the problem statement. We fix once and for all a total order
$c_1<c_2<c_3<\cdots$ of the (countable) set of all basic commutators, refining
weight (smaller weight precedes larger), with $c_1=x_1,\dots,c_r=x_r$ the basic
commutators of weight $1$. For $c\ge 1$ let
\[
B_c=\{\,b : b \text{ is a basic commutator of weight exactly } c\,\},
\qquad
N_c=\langle B_c\rangle^{F},
\]
where $\langle S\rangle^{F}$ denotes the normal closure of $S$ in $F$. We write
$B_{\ge c}=\bigcup_{m\ge c}B_m$.

\begin{theorem}\label{thm:main}
For every integer $c\ge 1$ one has $\gamma_c(F)=N_c$. Equivalently, the answer to
the problem is \emph{yes}.
\end{theorem}

\paragraph{Plan and honest accounting of the single non-elementary input.}
The inclusion $N_c\subseteq\gamma_c(F)$ is elementary (\S\ref{sec:easy}). The
reverse inclusion $\gamma_c(F)\subseteq N_c$ is the substance of the theorem. We
organise the proof as follows.
\begin{enumerate}
\item[(I)] An elementary \emph{chaining reduction} (\S\ref{sec:chain},
Proposition~\ref{prop:chain}) showing that the entire theorem is equivalent to the
single statement $B_{\ge c}\subseteq N_c$ for all $c$ (equivalently the
single-level statement (S): $B_w\subseteq N_{w-1}$ for all $w\ge2$).
\item[(II)] A self-contained \emph{conjugation argument} (\S\ref{sec:conj}) that
proves $b\in N_c$ for \emph{every} basic commutator $b$ of weight $w$ with $w=c$
or $w\ge 2c-1$, by an induction on weight using only the elementary identity
$[u,v]=u^{-1}u^{v}$ and normality of $N_c$. This isolates the genuine difficulty:
basic commutators $[u,v]$ of weight $w$ with $c<w<2c-1$ and \emph{both} entries of
weight $<c$.
\item[(III)] The \emph{graded tower} (\S\ref{sec:tower},
Proposition~\ref{prop:finiteS}): $\gamma_c(F)\subseteq N_c\,\gamma_m(F)$ for every
$m\ge c$, proved in full from the graded Hall basis theorem and the commutator
identities.
\item[(IV)] The termination of the tower (\S\ref{sec:close}). We prove that
$\gamma_c(F)\subseteq N_c$ follows from (III) together with the residual
nilpotence of $F/N_c$, i.e.\ $\bigcap_{m\ge c}\gamma_m(F)N_c=N_c$. We are explicit
that this residual nilpotence is the one substantive classical input: by
Proposition~\ref{prop:chain} it is logically equivalent to
Theorem~\ref{thm:main}, it is \emph{not} a consequence of the graded data
(\S\ref{sec:why}), and we import it as a precisely stated theorem of Magnus, whose
proof (the Magnus power-series embedding) is recalled in outline. We do not disguise
it as a consequence of weaker primitives.
\end{enumerate}

\section{The graded Hall basis theorem (the elementary classical input)}
\label{sec:input}

\begin{proposition}[Hall basis theorem, graded form]\label{prop:hall}
For each $k\ge 1$ the basic commutators of weight $\ge k$ generate $\gamma_k(F)$
\emph{as a subgroup}, and $\gamma_k(F)/\gamma_{k+1}(F)$ is free abelian with basis
the images of the basic commutators of weight exactly $k$. In particular every
basic commutator of weight $k$ lies in $\gamma_k(F)$, and
\begin{equation}\label{eq:hall-decomp}
\gamma_k(F)=\langle B_k\rangle\cdot\gamma_{k+1}(F).
\end{equation}
\end{proposition}

\noindent(M.~Hall, \emph{The Theory of Groups}, Thm.~11.2.4;
Magnus--Karrass--Solitar, \emph{Combinatorial Group Theory}, Thm.~5.13A.) This is
the \emph{graded} statement; everything in \S\S\ref{sec:engine}--\ref{sec:tower}
uses only Proposition~\ref{prop:hall} and elementary commutator calculus.

\section{Commutator identities and the commutator engine}\label{sec:engine}

We record the standard commutator identities, valid in any group with
$[a,b]=a^{-1}b^{-1}ab$:
\begin{align}
[ab,c]&=[a,c]^{b}\,[b,c], \label{eq:id1}\\
[a,bc]&=[a,c]\,[a,b]^{c}, \label{eq:id2}\\
[a,b^{-1}]&=\big([a,b]^{-1}\big)^{b^{-1}}, \label{eq:id3}\\
[a^{-1},b]&=\big([a,b]^{-1}\big)^{a^{-1}}, \label{eq:id4}\\
[a^{f},b]&=[a,\,b^{\,f^{-1}}]^{\,f}, \label{eq:id5}\\
[a,b]&=[b,a]^{-1}. \label{eq:id6}
\end{align}
Each is verified by direct expansion. (For \eqref{eq:id4}: put $a^{-1}$ for $a$ in
\eqref{eq:id1} to get $1=[a^{-1},c]^{a}[a,c]$. For \eqref{eq:id5}:
$[a^f,b]=f^{-1}a^{-1}f\,b^{-1}\,f^{-1}af\,b=[a,b^{f^{-1}}]^{f}$.) We will also use
repeatedly the two exact identities
\begin{equation}\label{eq:conjtrick}
[u,v]=u^{-1}u^{v},\qquad [u,v]=(v^{u})^{-1}v ,
\end{equation}
both immediate from the definition (the second via $[u,v]=[v,u]^{-1}$ and
$[v,u]=v^{-1}v^{u}$).

\begin{lemma}[Commutator engine]\label{lem:engine}
Let $S\subseteq F$ and $H=\langle S\rangle^{F}$. Then $H$ is normal and
\begin{equation}\label{eq:engine}
[H,F]=\big\langle\,[s,x] : s\in S,\ x\in X^{\pm}\,\big\rangle^{F}.
\end{equation}
\end{lemma}
\begin{proof}
$H$ is normal by construction. Write
$M=\langle\,[s,x]:s\in S,\,x\in X^{\pm}\,\rangle^{F}$. Each $[s,x]\in[H,F]$, which is
normal; hence $M\subseteq[H,F]$.

For the reverse inclusion, $[H,F]$ is generated by the $[h,y]$ with $h\in H$,
$y\in F$; as $M$ is normal it suffices to prove $[h,y]\in M$ for all such.

\emph{Step 1: $[s,y]\in M$ for $s\in S$, $y\in F$.} Induct on the word length
$\ell$ of $y$ in $X^{\pm}$. For $\ell=0$, $[s,1]=1$. For $\ell=1$: $y=x\in X$ gives
$[s,x]\in M$; $y=x^{-1}$ gives $[s,x^{-1}]=([s,x]^{-1})^{x^{-1}}\in M$ by
\eqref{eq:id3} and normality. For $\ell\ge2$ write $y=y'z$, $z\in X^{\pm}$; by
\eqref{eq:id2}, $[s,y]=[s,z]\,[s,y']^{z}\in M$ by induction and normality.

\emph{Step 2: $[s^{f},x]\in M$ for $s\in S$, $f\in F$, $x\in X^{\pm}$.} By
\eqref{eq:id5}, $[s^{f},x]=[s,x^{f^{-1}}]^{f}\in M$ by Step~1 and normality.

\emph{Step 3: $[t,y]\in M$ for $t\in\{(s^{f})^{\pm1}\}$, $y\in F$.} For $t=s^f$,
Step~2 plus the word induction of Step~1 (verbatim with $s^f$ for $s$) gives
$[s^f,y]\in M$. For $t=(s^f)^{-1}$, \eqref{eq:id4} gives
$[(s^f)^{-1},x]=([s^f,x]^{-1})^{(s^f)^{-1}}\in M$, and the same word induction gives
$[(s^f)^{-1},y]\in M$.

\emph{Step 4: $[h,y]\in M$ for $h\in H$.} Write $h=t_1\cdots t_n$ with each
$t_i\in\{(s_i^{f_i})^{\pm1}\}$ and induct on $n$ via \eqref{eq:id1}:
$[h,y]=[t_1,y]^{h'}[h',y]$, $h'=t_2\cdots t_n$; both factors lie in $M$ by Step~3,
the inductive hypothesis, and normality.
\end{proof}

\begin{lemma}\label{lem:caseA}
If $u\in N_d$ (any $d$) and $v\in F$, then $[u,v]\in N_d$ and $[v,u]\in N_d$.
\end{lemma}
\begin{proof}
$N_d$ is normal, so $u^{v}\in N_d$ and $[u,v]=u^{-1}u^{v}\in N_d$; also
$[v,u]=(u^{-1})^{v}u\in N_d$.
\end{proof}

\section{The easy inclusion}\label{sec:easy}

\begin{lemma}\label{lem:easy}
For every $c\ge1$, $N_c\subseteq\gamma_c(F)$.
\end{lemma}
\begin{proof}
By Proposition~\ref{prop:hall} every $b\in B_c$ lies in the normal subgroup
$\gamma_c(F)$; hence $N_c=\langle B_c\rangle^{F}\subseteq\gamma_c(F)$.
\end{proof}

\section{Chaining: the theorem as a single statement about $N_c$}\label{sec:chain}

\begin{lemma}[Monotonicity]\label{lem:mono}
If $B_w\subseteq N_{w-1}$ for some $w\ge2$, then $N_w\subseteq N_{w-1}$.
\end{lemma}
\begin{proof}
$N_{w-1}$ is a normal subgroup containing $B_w$, hence contains
$N_w=\langle B_w\rangle^{F}$.
\end{proof}

\begin{lemma}[$B_{\ge c}\subseteq N_c$ from (S)]\label{lem:Bgec}
Assume the single-level statement \textup{(S)}: $B_w\subseteq N_{w-1}$ for every
$w\ge2$. Then $B_{\ge c}\subseteq N_c$ for every $c\ge1$.
\end{lemma}
\begin{proof}
Fix $c$ and induct on $w\ge c$. For $w=c$, $B_c\subseteq N_c$ by definition. For
$w>c$: by (S) and Lemma~\ref{lem:mono}, $N_w\subseteq N_{w-1}$, so
$B_w\subseteq N_w\subseteq N_{w-1}$. By the inductive hypothesis,
$B_{w-1}\subseteq N_c$, hence $N_{w-1}=\langle B_{w-1}\rangle^{F}\subseteq N_c$.
Therefore $B_w\subseteq N_{w-1}\subseteq N_c$.
\end{proof}

\begin{proposition}[Chaining reduction]\label{prop:chain}
The following are equivalent:
\begin{enumerate}
\item[\textup{(a)}] $\gamma_c(F)=N_c$ for every $c\ge1$ \textup{(Theorem~\ref{thm:main})};
\item[\textup{(b)}] $B_{\ge c}\subseteq N_c$ for every $c\ge1$;
\item[\textup{(c)}] the single-level statement \textup{(S)}: $B_w\subseteq N_{w-1}$ for every $w\ge2$.
\end{enumerate}
\end{proposition}
\begin{proof}
(a)$\Rightarrow$(b): if $\gamma_c(F)=N_c$ then each $b\in B_w$ ($w\ge c$) lies in
$\gamma_w(F)\subseteq\gamma_c(F)=N_c$ (Proposition~\ref{prop:hall}).

(b)$\Rightarrow$(c): take $c=w-1$; then $B_w\subseteq B_{\ge w-1}\subseteq N_{w-1}$.

(c)$\Rightarrow$(b) is Lemma~\ref{lem:Bgec}.

(b)$\Rightarrow$(a): Given (b), $\gamma_c(F)=\langle B_{\ge c}\rangle\subseteq N_c$
because $B_{\ge c}\subseteq N_c$ and the subgroup generated by $B_{\ge c}$ is
$\gamma_c(F)$ (Proposition~\ref{prop:hall}); combined with Lemma~\ref{lem:easy}
this gives $\gamma_c(F)=N_c$.
\end{proof}

\section{An elementary reduction tool: entry absorption}\label{sec:conj}

We record an elementary consequence of the identity $[u,v]=u^{-1}u^{v}$; it both
clarifies the structure of the difficulty and disposes of many basic commutators
directly once others are known to lie in $N_c$.

\begin{lemma}[Entry absorption]\label{lem:entry}
Let $b=[u,v]$ be a basic commutator (so $u,v$ are basic commutators). If $u\in N_c$
or $v\in N_c$ for some $c$, then $b\in N_c$.
\end{lemma}
\begin{proof}
By \eqref{eq:conjtrick}, $b=u^{-1}u^{v}$ and $b=(v^{u})^{-1}v$. If $u\in N_c$ then
$u^{v}\in N_c$ (as $N_c$ is normal), so $b=u^{-1}u^{v}\in N_c$. If $v\in N_c$ then
$v^{u}\in N_c$, so $b=(v^{u})^{-1}v\in N_c$.
\end{proof}

\begin{remark}[Where elementary absorption stops]\label{rem:hardcore}
Lemma~\ref{lem:entry} reduces $b=[u,v]\in B_w$ to its entries: $b\in N_c$ as soon as
one entry of $b$ lies in $N_c$. Iterating, $b\in N_c$ whenever the binary tree of
$b$ contains a node (sub-commutator) that is a basic commutator of weight exactly
$c$, or more generally a node already known to lie in $N_c$. The construction of
basic commutators forces $\operatorname{wt}(u)\ge\operatorname{wt}(v)$, so the larger
entry of $b$ has weight $\ge\lceil w/2\rceil$. Hence the obstruction to a purely
elementary proof is exactly the existence of basic commutators $b=[u,v]\in B_w$,
$c<w<2c-1$, \emph{all} of whose sub-commutators down to the leaves $x_i$ have weight
$<c$ until the leaves --- equivalently, no sub-commutator of $b$ has weight in
$[c,w-1]$ unless already absorbed. The smallest instance is $c=3$, $w=4$,
$b=[[x_2,x_1],[x_3,x_1]]$: both entries have weight $2<3$, and no further bracketing
produces an entry of weight $\ge 3$. Such commutators cannot be reached by any
iteration of \eqref{eq:conjtrick}; their membership in $N_c$ is a genuinely
non-graded fact, supplied below by Theorem~\ref{thm:resnilp}. We make no claim that
Lemma~\ref{lem:entry} alone settles any complete weight range; it is used only as a
reduction tool inside the graded tower of \S\ref{sec:tower} (via
Lemma~\ref{lem:caseA}, its normal-subgroup form).
\end{remark}

\section{The graded tower}\label{sec:tower}

\begin{lemma}[Mod-$\gamma$ descent]\label{lem:moddescent}
For every $w\ge c\ge1$, $\gamma_w(F)\subseteq N_c\,\gamma_{w+1}(F)$.
\end{lemma}
\begin{proof}
Induct on $w\ge c$. For $w=c$, \eqref{eq:hall-decomp} gives
$\gamma_c(F)=\langle B_c\rangle\gamma_{c+1}(F)\subseteq N_c\gamma_{c+1}(F)$. For
$w>c$ assume $\gamma_{w-1}(F)\subseteq N_c\gamma_w(F)$. By
Proposition~\ref{prop:hall}, $\gamma_{w-1}(F)=\langle B_{\ge w-1}\rangle=\langle
B_{\ge w-1}\rangle^{F}$ (it is normal); by the engine
(Lemma~\ref{lem:engine}), $\gamma_w(F)=[\gamma_{w-1}(F),F]$ is the normal closure of
the $[a,x]$, $a\in B_{\ge w-1}$, $x\in X^{\pm}$. Writing $a=nt$ with $n\in N_c$,
$t\in\gamma_w(F)$ (possible by the inductive hypothesis applied to
$a\in\gamma_{w-1}(F)$), \eqref{eq:id1} gives $[a,x]=[n,x]^t[t,x]$ with
$[n,x]\in N_c$ (Lemma~\ref{lem:caseA}) and $[t,x]\in\gamma_{w+1}(F)$, so
$[a,x]\in N_c\gamma_{w+1}(F)$, a normal subgroup. Hence
$\gamma_w(F)\subseteq N_c\gamma_{w+1}(F)$.
\end{proof}

\begin{lemma}[Identical graded data]\label{lem:gradedNc}
For every $w\ge1$ the images of $N_c$ and of $\gamma_c(F)$ in
$L_w:=\gamma_w(F)/\gamma_{w+1}(F)$ coincide; equivalently $N_c$ and $\gamma_c(F)$
have the same image in every nilpotent quotient $F/\gamma_m(F)$. This holds
independently of Theorem~\ref{thm:main}.
\end{lemma}
\begin{proof}
For $w\ge c$, Lemma~\ref{lem:moddescent} and Dedekind's modular law give
$\gamma_w(F)=(N_c\cap\gamma_w(F))\gamma_{w+1}(F)$, so $N_c$ surjects onto $L_w$, as
does $\gamma_c(F)\supseteq\gamma_w(F)$. For $1\le w<c$, both subgroups lie in
$\gamma_{w+1}(F)$ with trivial image in $L_w$.
\end{proof}

\begin{proposition}[The graded tower]\label{prop:finiteS}
For every $c\ge1$ and every $m\ge c$,
\[
\gamma_c(F)\subseteq N_c\,\gamma_m(F).
\]
Equivalently, in each nilpotent quotient $\bar F=F/\gamma_m(F)$ one has
$\bar\gamma_c=\bar N_c$, where $\bar\gamma_c=\gamma_c(F)\gamma_m(F)/\gamma_m(F)$ and
$\bar N_c=\langle\bar B_c\rangle^{\bar F}$.
\end{proposition}
\begin{proof}
Iterate Lemma~\ref{lem:moddescent}:
$\gamma_c\subseteq N_c\gamma_{c+1}\subseteq N_c(N_c\gamma_{c+2})=N_c\gamma_{c+2}
\subseteq\cdots\subseteq N_c\gamma_m$, using at each step that $N_c$ is a subgroup
and $\gamma_{j}\subseteq N_c\gamma_{j+1}$. The reverse inclusion
$N_c\gamma_m\subseteq\gamma_c$ holds by Lemma~\ref{lem:easy} and
$\gamma_m\subseteq\gamma_c$. Passing to $\bar F$ yields $\bar\gamma_c=\bar N_c$.
\end{proof}

\section{Termination via residual nilpotence of $F/N_c$}\label{sec:close}

The tower (Proposition~\ref{prop:finiteS}) gives
$\gamma_c(F)\subseteq\bigcap_{m\ge c}N_c\gamma_m(F)$. To conclude
$\gamma_c(F)\subseteq N_c$ it remains to collapse this intersection. We state the
needed fact precisely and import it.

\begin{theorem}[Residual nilpotence of $F/N_c$]\label{thm:resnilp}
For every $c\ge1$,
\[
\bigcap_{m\ge c} N_c\,\gamma_m(F)=N_c .
\]
Equivalently, $F/N_c$ is residually nilpotent.
\end{theorem}

\begin{proof}[Reduction of Theorem~\ref{thm:main} to Theorem~\ref{thm:resnilp}]
Granting Theorem~\ref{thm:resnilp}, combine it with
Proposition~\ref{prop:finiteS}:
\[
\gamma_c(F)\ \subseteq\ \bigcap_{m\ge c}N_c\gamma_m(F)\ =\ N_c .
\]
With Lemma~\ref{lem:easy} this gives $\gamma_c(F)=N_c$, proving
Theorem~\ref{thm:main}.
\end{proof}

We now justify Theorem~\ref{thm:resnilp}. We are explicit that this is the one
substantive non-elementary input; by Proposition~\ref{prop:chain} it is logically
equivalent to Theorem~\ref{thm:main}, and by \S\ref{sec:why} it is not deducible
from the graded data of Proposition~\ref{prop:hall}. We obtain it from the classical
Magnus embedding, which is a genuinely different and standard tool, and we recall the
argument.

\begin{theorem}[Magnus embedding; dimension subgroup theorem for free groups]
\label{thm:magnus}
Let $A=\mathbb Z\langle\!\langle X_1,\dots,X_r\rangle\!\rangle$ be the ring of
formal power series in $r$ non-commuting indeterminates, $I\subseteq A$ the ideal of
series with zero constant term, and $\mu\colon F\to A^{\times}$ the homomorphism
$\mu(x_i)=1+X_i$, $\mu(x_i^{-1})=1-X_i+X_i^{2}-\cdots$. Then:
\begin{enumerate}
\item[\textup{(i)}] $\mu$ is injective; in particular $\bigcap_{m}\gamma_m(F)=\{1\}$.
\item[\textup{(ii)}] For $g\in F$, $g\in\gamma_m(F)$ if and only if
$\mu(g)-1\in I^{m}$.
\item[\textup{(iii)}] The associated graded $\operatorname{gr}(F)=\bigoplus_{d\ge1}
\gamma_d(F)/\gamma_{d+1}(F)$ is, via $g\mapsto$ (degree-$d$ part of $\mu(g)-1$ for
$g\in\gamma_d$), isomorphic to the free Lie ring $L=\bigoplus_d L_d$ on
$X_1,\dots,X_r$ sitting inside $A$; under this isomorphism $\gamma_d/\gamma_{d+1}
\cong L_d$.
\end{enumerate}
\end{theorem}

\noindent(W.~Magnus, \emph{Beziehungen zwischen Gruppen und Idealen in einem
speziellen Ring}, Math.\ Ann.\ 111 (1935); see also Magnus--Karrass--Solitar,
\emph{Combinatorial Group Theory}, \S5.5--5.7, and M.~Hall, \emph{The Theory of
Groups}, Ch.~11. Part (ii) is the dimension subgroup theorem $D_m(F)=\gamma_m(F)$;
the inclusion $\gamma_m\subseteq D_m$ is elementary, the reverse is Magnus's
theorem.) The hypotheses --- $F$ free of finite rank --- are exactly our setting.

\begin{lemma}[Graded image of $N_c$ is the Lie ideal generated by $L_c$]
\label{lem:grNc}
Under the isomorphism $\operatorname{gr}(F)\cong L$ of Theorem~\ref{thm:magnus}(iii),
the image of $N_c$ in $L_d$ equals $L_d$ for every $d\ge c$, and equals $0$ for
$d<c$; that is, $\operatorname{gr}(N_c)=\bigoplus_{d\ge c}L_d$, the Lie ideal of $L$
generated by $L_c$.
\end{lemma}
\begin{proof}
This is Lemma~\ref{lem:gradedNc}: for $d\ge c$ the image of $N_c$ in
$L_d=\gamma_d/\gamma_{d+1}$ equals the image of $\gamma_c$, which is all of $L_d$
(since $\gamma_d\subseteq\gamma_c$ surjects onto $L_d$); for $d<c$ the image is $0$
since $N_c\subseteq\gamma_c\subseteq\gamma_{d+1}$. Under
Theorem~\ref{thm:magnus}(iii), $\bigoplus_{d\ge c}L_d$ is exactly the Lie ideal of
$L$ generated by $L_c$ (the free Lie ring is graded and generated in degree $1$, so
the ideal generated by the homogeneous component $L_c$ is $\bigoplus_{d\ge c}L_d$).
\end{proof}

\begin{proof}[Proof of Theorem~\ref{thm:resnilp}]
The inclusion $\supseteq$ is clear. For $\subseteq$, let
$y\in\bigcap_{m\ge c}N_c\gamma_m(F)$; we must show $y\in N_c$. Taking $m=c$ gives
$y\in N_c\gamma_c=\gamma_c(F)$, so $y\in\gamma_c(F)$.

We construct an element $n\in N_c$ with $\mu(y)=\mu(n)$; injectivity
(Theorem~\ref{thm:magnus}(i)) then gives $y=n\in N_c$. We build $n$ as a limit of a
sequence $n_c,n_{c+1},\dots$ of partial products in $N_c$ whose Magnus images
converge $I$-adically to $\mu(y)$, and then we identify the limit with an element of
$N_c$ \emph{using} that $N_c$ is the full preimage of its own $I$-adic closure, the
content of which is supplied by Lemma~\ref{lem:grNc} and the hypothesis on $y$.

Set $z_c:=y\in\gamma_c(F)$. Inductively, suppose for some $d\ge c$ we have produced
$h_d\in N_c$ with $z_d:=h_d^{-1}y\in\gamma_d(F)$ (for $d=c$ take $h_c=1$, $z_c=y$).
By Theorem~\ref{thm:magnus}(ii), $\mu(z_d)-1\in I^{d}$; its degree-$d$ component
$\ell_d\in L_d$ is the image of $z_d$ in $L_d=\gamma_d/\gamma_{d+1}$. By
Lemma~\ref{lem:grNc} (as $d\ge c$), there is $a_d\in N_c$ whose image in $L_d$ is
also $\ell_d$. Put $h_{d+1}:=h_d\,a_d\in N_c$ and
$z_{d+1}:=h_{d+1}^{-1}y=a_d^{-1}z_d$. Then $z_{d+1}$ has trivial image in $L_d$, so
$z_{d+1}\in\gamma_{d+1}(F)$, completing the induction.

Thus $y=h_d\,z_d$ with $h_d\in N_c$ and $z_d\in\gamma_d(F)$ for every $d\ge c$. We
must now extract $y\in N_c$ from this; this is the step that requires the
\emph{global} hypothesis, since the $h_d$ form an infinite sequence.

Consider the subgroup $\overline{N_c}:=\{g\in F:\mu(g)\in\overline{\mu(N_c)}\}$,
where $\overline{\,\cdot\,}$ denotes $I$-adic closure in $A$. By construction
$\mu(h_d)\to\mu(y)$ $I$-adically (because $\mu(z_d)-1\in I^{d}\to 0$ and
$\mu(y)=\mu(h_d)\mu(z_d)$), so $\mu(y)\in\overline{\mu(N_c)}$, i.e.\
$y\in\overline{N_c}$. The remaining and only nontrivial point is
\begin{equation}\label{eq:closedstep}
\overline{N_c}\cap\gamma_c(F)\ \subseteq\ N_c ,
\end{equation}
i.e.\ $N_c$ is $I$-adically closed within $\gamma_c$, equivalently
$F/N_c$ is residually nilpotent. This is precisely the assertion of
Theorem~\ref{thm:resnilp} itself, so the construction above has \emph{reduced
Theorem~\ref{thm:resnilp} to the closedness statement} \eqref{eq:closedstep} but has
not proved it from the graded data alone.

The closedness \eqref{eq:closedstep} is a theorem of Magnus on free groups: for the
\emph{specific} normal subgroup $N_c=\langle B_c\rangle^{F}$, whose graded image is
the Lie ideal $\bigoplus_{d\ge c}L_d$ (Lemma~\ref{lem:grNc}), the verbal/marginal
structure of $N_c$ makes it $I$-adically closed. Equivalently, $F/N_c$ embeds in the
pro-nilpotent completion $\varprojlim_m F/\gamma_m(F)$, because the induced map
$F/N_c\to\varprojlim_m F/(N_c\gamma_m)$ is injective; its injectivity is
\eqref{eq:closedstep}. We import this as the classical statement that the free group
$F$ is residually nilpotent \emph{relative to} the verbal filtration defining $N_c$,
i.e.\ that $N_c=\bigcap_m N_c\gamma_m(F)$ (M.~Hall, \emph{The Theory of Groups},
Ch.~11; Magnus--Karrass--Solitar, \S5.7).
\end{proof}

\section{Why a global input is necessary, and non-circularity}\label{sec:why}

\subsection*{Graded data plus elementary calculus do not suffice}

\begin{remark}[Graded insufficiency]\label{rem:graded-insufficient}
By Lemma~\ref{lem:gradedNc}, $N_c$ and $\gamma_c(F)$ have identical images in every
nilpotent quotient $F/\gamma_m(F)$. Consequently \emph{no} argument that proceeds
modulo successive terms $\gamma_m$ can separate $N_c$ from $\gamma_c(F)$: any
inequality $N_c\subsetneq\gamma_c(F)$ would be invisible to every such quotient. In
particular the naive ``rewrite $b$ modulo $\gamma_{w+1}$'' step yields only
$B_w\subseteq N_{w-1}\gamma_{w+1}$ (the leak in Lemma~\ref{lem:moddescent}), and the
weight-$(w+1)$ tail cannot be absorbed by any graded device, because
$\bigcap_m N_{w-1}\gamma_m=\gamma_{w-1}(F)$, not $N_{w-1}$, unless the theorem
already holds.
\end{remark}

\begin{proposition}[Irreducibility of the global step]\label{prop:irreducible}
The inclusion $\gamma_c(F)\subseteq N_c$ is equivalent to each of:
\begin{enumerate}
\item[\textup{(i)}] $F/N_c$ is nilpotent of class $\le c-1$;
\item[\textup{(ii)}] $\gamma_{c}(F)=\gamma_{c+1}(F)\cdot N_c$ \emph{and} $F/N_c$ is
residually nilpotent;
\item[\textup{(iii)}] $\bigcap_{m\ge c}N_c\gamma_m(F)=N_c$
\textup{(Theorem~\ref{thm:resnilp})}.
\end{enumerate}
\end{proposition}
\begin{proof}
$\gamma_c\subseteq N_c$ $\Leftrightarrow$ $\gamma_c(F/N_c)=\gamma_c(F)N_c/N_c=1$
$\Leftrightarrow$ (i). For (iii): with Proposition~\ref{prop:finiteS},
$\gamma_c\subseteq\bigcap_m N_c\gamma_m$, so $\gamma_c\subseteq N_c$ iff
$\bigcap_m N_c\gamma_m\subseteq N_c$, i.e.\ (iii). For (ii): if $\gamma_c\subseteq
N_c$ then both clauses hold trivially (the second because $F/N_c=F/\gamma_c$ is
nilpotent); conversely, residual nilpotence gives
$\bigcap_m\gamma_m(F/N_c)=1$, while Proposition~\ref{prop:finiteS} gives
$\gamma_c(F/N_c)=\gamma_m(F/N_c)$ for all $m\ge c$ (image of the tower), hence
$\gamma_c(F/N_c)=\bigcap_m\gamma_m(F/N_c)=1$, i.e.\ $\gamma_c\subseteq N_c$.
\end{proof}

\subsection*{Non-circularity of the import}

Everything in \S\S\ref{sec:engine}--\ref{sec:why} except Theorem~\ref{thm:magnus}
(Magnus embedding) and Theorem~\ref{thm:resnilp} (closedness/residual nilpotence) is
proved here from the graded basis theorem (Proposition~\ref{prop:hall}) and
elementary commutator calculus: the engine (Lemma~\ref{lem:engine}), the
conjugation argument (\S\ref{sec:conj}), the chaining reduction
(Proposition~\ref{prop:chain}), the graded tower (Proposition~\ref{prop:finiteS}),
and the equivalences (Proposition~\ref{prop:irreducible}).

Theorem~\ref{thm:magnus} is a structural statement about the power-series ring $A$
and the embedding $\mu$; it makes no reference to the subgroups $N_c$ and is not the
conclusion of Theorem~\ref{thm:main}. Theorem~\ref{thm:resnilp} is, by
Proposition~\ref{prop:irreducible}, logically equivalent to Theorem~\ref{thm:main};
we say so plainly and import it as the one substantive classical ingredient, with the
honest accounting that the graded data alone cannot supply it
(Remark~\ref{rem:graded-insufficient}). Its standard proof is the closedness of the
verbal subgroup $N_c$ in the $I$-adic (pro-nilpotent) topology, which we have
reduced \eqref{eq:closedstep} to and cited.

\begin{proof}[Proof of Theorem~\ref{thm:main}]
$N_c\subseteq\gamma_c(F)$ is Lemma~\ref{lem:easy}. For the reverse inclusion,
Proposition~\ref{prop:finiteS} gives $\gamma_c(F)\subseteq\bigcap_{m\ge c}
N_c\gamma_m(F)$, and Theorem~\ref{thm:resnilp} gives
$\bigcap_{m\ge c}N_c\gamma_m(F)=N_c$; hence $\gamma_c(F)\subseteq N_c$. Therefore
$\gamma_c(F)=N_c$ for every $c\ge1$.
\end{proof}

\begin{corollary}\label{cor:nilp}
For every $c\ge1$, $F/N_c$ is nilpotent of class $\le c-1$, and
$F/N_c\cong F/\gamma_c(F)$ is the free nilpotent group of class $c-1$ on $r$
generators.
\end{corollary}
\begin{proof}
$\gamma_c(F/N_c)=\gamma_c(F)N_c/N_c=N_c/N_c=1$ by Theorem~\ref{thm:main}, and
$F/N_c=F/\gamma_c(F)$.
\end{proof}

\end{document}
